\lecture{11}
\title{Linear Programming (Part I)}

\begin{document}

\textbf{Definition (Linear Programming).} We seek to optimize a linear objective function subject to linearity equality and linear inequality constraints.  It will always look like:
\[ \text{Max / min of } \left(\sum_{i = 1}^n c_i x_i\right) \text{ given } x_i \in \mathbb{R}^{\geq 0} \text{ and } \left \{ \left(\sum_{i = 1}^n A_{j, i} x_i \ \right) \bigstar_j \ b_j \right \}_{j = 1}^m \text{ for } \bigstar_j \in \{ \leq , \geq , =\}. \]

\begin{problem}{Claw Machine}
    You have two claw machines $K_1$ and $K_2$.  Putting $\$x_1$ into $K_1$ yields a return of $\$ \left(\frac{x_1}{3}\right)$, and putting $\$x_2$ into $K_2$ yields a return of $\$ \left(\frac{x_2}{2}\right)$.  Also, you must put at least as much into $K_1$ as $K_2$.  You have $\$60$.
\end{problem}

\begin{solution*}{Claw Machine LP}
    We can write the above setup as a linear program like so:
    \[ \text{Max / min of } \dfrac{1}{3}x_1 + \dfrac{1}{2}x_2 \text{ given } x_i \in \mathbb{R}^{\geq 0} \text{ and } x_1 - x_2 \geq 0 \text{ and } x_1 + x_2 \leq 60. \]
\end{solution*}

We can determine a \emph{normal form} for linear programs.
\begin{itemize}
    \item Finding the max of $\mathcal{P}$ is the same as finding the min of $-\mathcal{P}$.  So we'll default max.
    \item If $\bigstar_j$ is only allowed to be $\leq$, then we can also implement constraints with $\geq$ and $=$.
    \begin{itemize}
        \item (Using \emph{slack variables}, we can implement $\{\leq, \geq, =\}$ using just $=$.  But we'll use $\leq$.)
    \end{itemize}
\end{itemize}
We can also rewrite our constraints in linear algebra notation.  (Here, we say $\vec{m} \leq \vec{n}$ iff $m_i \leq n_i$ for all indices $i$.)
\[ \text{Find max of } \vec{c} \cdot \vec{x} \text{ given } \vec{x} \geq \vec{0} \text{ and } A\vec{x} \leq \vec{b}. \]
\textbf{Definition.} Given a set of constraints given by $A \vec{x} \leq \vec{b}$ of a linear program $\mathcal{P}$,
\begin{itemize}
    \item We say a \textbf{feasible solution} is some $\vec{\tilde{x}}$ such that $A\vec{\tilde{x}} \leq \vec{b}$.
    \item We say $\textsc{Feasible}(\mathcal{P})$ is the set of all $\vec{x} \in \mathbb{R}^{\geq 0}$ such that $A\vec{x} \leq \vec{b}$.
    \item Some constraints might be \textbf{degenerate} in that they can be removed without changing $\textsc{Feasible}$.
    \begin{itemize}
        \item This can happen by \emph{weakening} via $\{x \geq 2, x \geq 0\} \to \{x \geq 2\}$.
        \item This can happen by \emph{linear combination} via $\{a \geq 2, b \geq 3, a + b \geq 5\} \to \{a \geq 2, b \geq 3\}$.
    \end{itemize}
    \item We say the \textbf{optimal value}, possibly $\infty$, is achieved by some \textbf{optimal solution}.
    \begin{itemize}
        \item If the optimal value is finite, then the LP is \textbf{bounded}.  Otherwise, the LP is \textbf{unbounded}.
        \item If $\textsc{Feasible}(\mathcal{P})$ is bounded (topologically speaking), then LP must be bounded.
        \item There is always a unique optimal value, possibly $\infty$.  \emph{(Pandora's box\dots)}
        \item There may not always be a unique optimal solution, though.
        \begin{itemize}
            \item \emph{Exercise:} Prove that if there are two optimal solutions, then there are infinitely many.
        \end{itemize}
    \end{itemize}
\end{itemize}
The optimal solution always occurs at a corner point, obviously.  There exist polynomial-time solving algorithms.

There's this idea of duality that's also pretty obvious.  You want to \emph{minimize} the RHS of a linear combination $\vec{y}$ of the constraints $A \vec{x} \leq \vec{b}$, constrained on the LHS of this linear combination being at most $\vec{c} \cdot \vec{x}$.

\end{document}